% =====================================================================
%  บทที่ 1  บทนำ
% =====================================================================
\chapter{บทนำ}
\label{ch:intro}

\section{ความเป็นมาและความสำคัญของปัญหา}

ปัญญาประดิษฐ์เชิงสร้างสรรค์ (Generative Artificial Intelligence) เป็นเทคโนโลยีที่สามารถสร้างเนื้อหาใหม่ เช่น ข้อความ ภาพ เสียง และโค้ดโปรแกรม ได้จากข้อมูลที่ใช้ฝึกสอน ความก้าวหน้าของสถาปัตยกรรมทรานส์ฟอร์เมอร์ (Transformer)~\cite{vaswani2017attention} ทำให้เกิดแบบจำลองภาษาขนาดใหญ่ที่สามารถเรียนรู้จากตัวอย่างเพียงไม่กี่ตัวอย่างได้~\cite{brown2020language} ส่งผลให้มีการนำเทคโนโลยีนี้มาประยุกต์ใช้ในด้านการศึกษาอย่างแพร่หลาย

ในรายวิชาการเขียนโปรแกรมเบื้องต้น ผู้เรียนจำนวนมากประสบปัญหาในการทำความเข้าใจข้อความแจ้งข้อผิดพลาด (error message) และไม่สามารถแก้ไขโค้ดได้ด้วยตนเอง ขณะที่ผู้สอนมีเวลาจำกัดในการให้คำแนะนำแก่ผู้เรียนเป็นรายบุคคล\footnote{ตัวอย่างเชิงอรรถ: ข้อมูลจากการสำรวจเบื้องต้นของผู้วิจัยในภาคการศึกษาที่ 1/2568} ด้วยเหตุนี้ ผู้วิจัยจึงสนใจพัฒนาระบบผู้ช่วยสอนที่สามารถให้คำแนะนำแก่ผู้เรียนได้ตลอดเวลา

\section{วัตถุประสงค์ของการวิจัย}

\begin{enumerate}
  \item เพื่อพัฒนาระบบผู้ช่วยสอนด้วยปัญญาประดิษฐ์เชิงสร้างสรรค์สำหรับรายวิชาการเขียนโปรแกรมเบื้องต้น
  \item เพื่อเปรียบเทียบผลสัมฤทธิ์ทางการเรียนของผู้เรียนที่ใช้และไม่ใช้ระบบผู้ช่วยสอน
  \item เพื่อศึกษาความพึงพอใจของผู้เรียนที่มีต่อระบบผู้ช่วยสอน
\end{enumerate}

\section{ขอบเขตของการวิจัย}

\subsection{ขอบเขตด้านเนื้อหา}
ระบบครอบคลุมเนื้อหาการเขียนโปรแกรมภาษา Python ได้แก่ ตัวแปร เงื่อนไข การวนซ้ำ ฟังก์ชัน และรายการ (list)

\subsection{ขอบเขตด้านประชากรและกลุ่มตัวอย่าง}
ประชากรคือนักศึกษาชั้นปีที่ 1 ที่ลงทะเบียนเรียนรายวิชาการเขียนโปรแกรมเบื้องต้น ภาคการศึกษาที่ 1 ปีการศึกษา 2569 กลุ่มตัวอย่างจำนวน 60 คน ได้มาโดยการสุ่มแบบกลุ่ม (cluster random sampling)

\section{ประโยชน์ที่คาดว่าจะได้รับ}

\begin{itemize}
  \item ได้ระบบผู้ช่วยสอนที่ผู้เรียนสามารถใช้งานได้ตลอด 24 ชั่วโมง
  \item ผู้สอนมีเวลามากขึ้นในการดูแลผู้เรียนที่ต้องการความช่วยเหลือเป็นพิเศษ
  \item เป็นแนวทางในการประยุกต์ใช้ปัญญาประดิษฐ์เชิงสร้างสรรค์ในรายวิชาอื่น
\end{itemize}

\section{นิยามศัพท์เฉพาะ}

\textbf{ปัญญาประดิษฐ์เชิงสร้างสรรค์} หมายถึง ระบบปัญญาประดิษฐ์ที่สามารถสร้างข้อมูลใหม่ที่มีลักษณะคล้ายกับข้อมูลที่ใช้ฝึกสอน

\textbf{ผู้ช่วยสอน} หมายถึง ระบบซอฟต์แวร์ที่ตอบคำถามและให้คำแนะนำแก่ผู้เรียนในรายวิชา โดยอ้างอิงจากเอกสารประกอบการสอน
